\section{Voting Rights Act Compliance}
\label{sec:vra}

\subsection{The Callais Disentanglement Problem}

Any analysis of algorithmic redistricting and the Voting Rights Act must begin with the core tension identified in \textit{Louisiana v.\ Callais} (2025) and the line of cases addressing the interaction between VRA compliance and partisan gerrymandering.\footnote{See \textit{Louisiana v.\ Callais}, --- U.S. --- (2025); \textit{Allen v.\ Milligan}, 599 U.S. 1 (2023); \textit{Alexander v.\ South Carolina State Conference of the NAACP}, 144 S. Ct. 1221 (2024).} The problem is epistemic: when a majority-minority district is drawn, it is often impossible to determine whether the district's shape was driven by race (potentially required by VRA Section 2) or by partisanship (potentially impermissible under the Constitution's equal protection guarantee, see \textit{Shaw v.\ Reno}).\footnote{\textit{Shaw v.\ Reno}, 509 U.S. 630 (1993).} Because racial and partisan voting patterns are correlated in much of the country, any algorithm that draws majority-minority districts will also affect partisan compositions.

Algorithmic redistricting resolves this disentanglement problem structurally. The baseline algorithm uses no racial or partisan data; it optimizes for population equality and compactness alone. Any racial or partisan pattern in its outputs reflects the interaction between geographic voter distribution and compactness optimization, not a conscious choice to favor any racial or partisan group. When the baseline algorithm's outputs are subjected to a VRA analysis---examining whether minority populations are packed into any single district to dilute their influence, or cracked across multiple districts to prevent any majority-minority district from forming---the analysis examines what geographic compactness alone produces, providing a clean baseline against which VRA adjustment can be measured.

The VRA mode of the algorithm then operates as a documented, transparent modification: it increases the weight on edges connecting census tracts with high minority populations, making it more expensive for the partitioner to cut through minority communities. This edge-weighting approach, described formally in \citet{dellaLibera2026vra}, creates majority-minority districts by preserving the geographic clustering of minority populations---it does not ``target'' any specific percentage of minority voters but instead prevents the algorithm from fragmenting natural geographic concentrations. Because the modification is mathematical and documented, courts can examine precisely what the edge-weighting did and did not do. The disentanglement problem that plagued \textit{Callais} does not arise because the record contains a baseline (no VRA mode), a modification (VRA-weighted mode), and a documented explanation of exactly what data drove each.

\subsection{Section 2 and the \textit{Gingles} Framework}

Section 2 of the Voting Rights Act prohibits voting practices or procedures that result in the denial or abridgement of the right to vote on account of race, color, or membership in a language minority group.\footnote{52 U.S.C. \S\,10301.} Under \textit{Thornburg v.\ Gingles} (1986), a Section 2 vote dilution claim based on single-member district lines requires proof of three preconditions:\footnote{\textit{Thornburg v.\ Gingles}, 478 U.S. 30, 50-51 (1986).} (1) the minority group is ``sufficiently large and geographically compact to constitute a majority'' in a single-member district; (2) the minority group is ``politically cohesive''; and (3) the majority votes ``sufficiently as a bloc to enable it ... usually to defeat the minority's preferred candidate.''

Algorithmic redistricting operates cleanly within this framework. Prong 1 (geographic compactness) is directly served by the algorithm's compactness optimization: the algorithm is less likely than a human mapmaker to create artificial geographic configurations that spread a geographically compact minority population across multiple districts. Empirical evidence supports this: the algorithm analyzed in \citet{dellaLibera2026vra} produces 69 more majority-minority districts than the enacted plans for the same states, demonstrating that compactness-optimizing redistricting serves rather than undermines Section 2's geographic-compactness requirement.

Prongs 2 and 3 (political cohesion and majority bloc voting) are not affected by the algorithm at all---they are empirical questions about voter behavior that neither the algorithm nor the redistricting process determines. If a community of color is politically cohesive and faces majority bloc voting sufficient to defeat its preferred candidates, these facts exist in the world independent of how district lines are drawn. The algorithm's contribution is to ensure that district lines do not artificially prevent the satisfaction of prong 1 through non-compact configurations that dilute geographically concentrated minority populations.

\subsection{The VRA Safe Harbor}

The model statute includes a safe harbor provision for majority-minority districts produced by the VRA mode of the algorithm. States that run the algorithm in VRA mode, document the edge-weighting parameters, and produce districts that satisfy the \textit{Gingles} preconditions for identified minority communities shall be presumed to have satisfied Section 2, with the burden shifting to any challenger to demonstrate that an alternative configuration would produce additional Section 2 districts without violating any other statutory criterion.

This safe harbor addresses two concerns simultaneously. It protects states from Section 2 liability when the algorithm's VRA mode produces majority-minority districts consistent with geographic concentration of minority populations. And it protects those majority-minority districts from \textit{Shaw v.\ Reno} racial gerrymandering challenges by establishing that the district's shape was driven by documented mathematical parameters applied to geographic clustering data, not by a conscious effort to maximize any particular racial percentage.

The \textit{Shaw} doctrine's concern is with districts whose ``bizarre'' shapes can be explained only by racial motivation.\footnote{\textit{Shaw v.\ Reno}, 509 U.S. at 644.} Districts produced by a compactness-optimizing algorithm with documented VRA-mode edge-weighting have a complete non-racial explanation for every aspect of their shape: compactness optimization explains the general configuration, and geographic minority-population clustering explains why the VRA-mode adjustment preserved particular boundaries. The racial motivation concern simply does not apply to a process that can explain every boundary decision in terms of documented, non-racial optimization parameters.

\subsection{Post-\textit{Allen v.\ Milligan} Landscape}

\textit{Allen v.\ Milligan} (2023) reaffirmed the Section 2 framework and rejected Alabama's argument that race-conscious redistricting is unconstitutional whenever an alternative race-neutral plan would produce fewer majority-minority districts.\footnote{\textit{Allen v.\ Milligan}, 599 U.S. 1 (2023).} The Court held that Section 2 sometimes requires majority-minority districts and that the state must demonstrate that its plan does not result in minority vote dilution under the \textit{Gingles} framework.

Algorithmic redistricting is well-positioned under \textit{Allen}. Because the algorithm's baseline mode produces no majority-minority districts through any intentional design, any majority-minority districts in the baseline output reflect natural geographic concentration of minority populations---they satisfy the \textit{Gingles} compactness prong by construction. The VRA mode produces additional majority-minority districts where the baseline might have fragmented geographically concentrated populations. The net result, documented in \citet{dellaLibera2026vra}, is that the algorithm produces more majority-minority districts than enacted plans without ever targeting any specific racial percentage, precisely because it does not fragment naturally compact minority communities.

\textit{Alexander v.\ South Carolina} (2024) added a new wrinkle by raising the evidentiary standard for proving racial gerrymandering in the context of districts where race and partisanship are correlated.\footnote{\textit{Alexander v.\ South Carolina State Conference of the NAACP}, 144 S. Ct. 1221 (2024).} Algorithmic redistricting addresses this concern directly: the algorithm produces a documented record of exactly what data influenced each boundary decision. Where race and partisanship are correlated, the algorithm's record can demonstrate which factor---if either---drove any particular boundary. In most cases, the answer will be neither: the boundary will be driven by the compactness optimization of geographic adjacency, which is independent of both race and partisanship.
